\documentclass[11pt,a4paper]{article}

% Packages
\usepackage[utf8]{inputenc}
\usepackage[T1]{fontenc}
\usepackage{geometry}
\usepackage{listings}
\usepackage{xcolor}
\usepackage{booktabs}
\usepackage{hyperref}
\usepackage{fancyhdr}
\usepackage{tcolorbox}
\usepackage{enumitem}

% Page geometry
\geometry{margin=1in}

% Colors
\definecolor{codegreen}{rgb}{0,0.6,0}
\definecolor{codegray}{rgb}{0.5,0.5,0.5}
\definecolor{backcolour}{rgb}{0.95,0.95,0.92}
\definecolor{infoблue}{rgb}{0.9,0.95,1}

% Code listings style
\lstdefinestyle{bashstyle}{
    backgroundcolor=\color{backcolour},
    basicstyle=\ttfamily\footnotesize,
    breaklines=true,
    frame=single,
    framerule=0.5pt,
}
\lstset{style=bashstyle}

% Header/Footer
\pagestyle{fancy}
\fancyhf{}
\rhead{PIQC Tool Testing}
\lhead{GCP Permissions Request}
\rfoot{Page \thepage}

% Hyperref setup
\hypersetup{
    colorlinks=true,
    linkcolor=blue,
    urlcolor=blue,
}

\begin{document}

% Title
\begin{center}
    {\LARGE\textbf{GCP Permissions Request}}\\[0.5em]
    {\Large For PIQC Tool Testing}\\[1em]
    \textit{Date: January 11, 2026}
\end{center}

\vspace{1em}

%==============================================================================
\section{Overview}
%==============================================================================

We need to test the \textbf{PIQC (Production Inference Quality Control)} tool on your GCP environment. This tool discovers and documents AI/ML inference deployments running on Kubernetes.

To complete the testing, we need certain GCP permissions or your assistance in setting up the test environment.

%==============================================================================
\section{Account Information}
%==============================================================================

\begin{tcolorbox}[colback=backcolour,colframe=black!50,title=\textbf{Account to Grant Access}]
\textbf{Email:} \texttt{ammarkhan2264388@gmail.com}
\end{tcolorbox}

%==============================================================================
\section{Option 1: Grant Editor Access (Recommended)}
%==============================================================================

The simplest approach is to grant \textbf{Editor} role to my account. This allows me to:
\begin{itemize}
    \item Enable required APIs
    \item Create a test GKE cluster
    \item Deploy test workloads
    \item Run the PIQC tool
    \item Delete resources after testing
\end{itemize}

\subsection{How to Grant Editor Access}

\begin{enumerate}
    \item Go to \textbf{IAM \& Admin} page:\\
    \url{https://console.cloud.google.com/iam-admin/iam}
    
    \item Make sure you have the correct project selected in the top dropdown
    
    \item Click \textbf{"+ GRANT ACCESS"} button at the top
    
    \item In the \textbf{"New principals"} field, enter:\\
    \texttt{ammarkhan2264388@gmail.com}
    
    \item In the \textbf{"Select a role"} dropdown, choose:\\
    \texttt{Project > Editor}
    
    \item Click \textbf{"Save"}
\end{enumerate}

\begin{tcolorbox}[colback=backcolour,colframe=green!50!black,title=\textbf{After Granting Access}]
Please notify me once access is granted, and provide the \textbf{Project ID} so I can proceed with testing.
\end{tcolorbox}

%==============================================================================
\section{Option 2: Minimal Permissions (More Secure)}
%==============================================================================

If you prefer not to grant full Editor access, please grant these specific roles:

\begin{table}[h]
\centering
\begin{tabular}{@{}ll@{}}
\toprule
\textbf{Role} & \textbf{Purpose} \\
\midrule
Kubernetes Engine Admin & Create/manage GKE clusters \\
Compute Admin & Manage compute resources \\
Service Usage Admin & Enable APIs \\
\bottomrule
\end{tabular}
\caption{Minimum required roles for testing}
\end{table}

\subsection{How to Grant Multiple Roles}

\begin{enumerate}
    \item Go to: \url{https://console.cloud.google.com/iam-admin/iam}
    \item Click \textbf{"+ GRANT ACCESS"}
    \item Enter: \texttt{ammarkhan2264388@gmail.com}
    \item Click \textbf{"+ ADD ANOTHER ROLE"} for each role
    \item Add all three roles listed above
    \item Click \textbf{"Save"}
\end{enumerate}

%==============================================================================
\section{Option 3: You Set Up the Environment}
%==============================================================================

If you prefer to set up the environment yourself, please run the following commands in Google Cloud Shell (\url{https://shell.cloud.google.com}):

\subsection{Step 1: Enable Required APIs}

\begin{lstlisting}[language=bash]
gcloud services enable container.googleapis.com
gcloud services enable compute.googleapis.com
\end{lstlisting}

\subsection{Step 2: Create GKE Cluster with GPU}

\begin{lstlisting}[language=bash]
gcloud container clusters create piqc-test \
  --zone=us-central1-a \
  --machine-type=n1-standard-4 \
  --accelerator=type=nvidia-tesla-t4,count=1 \
  --num-nodes=1 \
  --spot \
  --disk-size=50GB
\end{lstlisting}

\subsection{Step 3: Install GPU Drivers}

\begin{lstlisting}[language=bash]
kubectl apply -f https://raw.githubusercontent.com/GoogleCloudPlatform/container-engine-accelerators/master/nvidia-driver-installer/cos/daemonset-preloaded.yaml
\end{lstlisting}

\subsection{Step 4: Share Cluster Access with Me}

\begin{lstlisting}[language=bash]
# Export kubeconfig file
kubectl config view --minify --flatten > piqc-kubeconfig.yaml
\end{lstlisting}

Send me the \texttt{piqc-kubeconfig.yaml} file securely, and I can test the PIQC tool from my machine.

%==============================================================================
\section{Estimated Costs}
%==============================================================================

\begin{table}[h]
\centering
\begin{tabular}{@{}lll@{}}
\toprule
\textbf{Resource} & \textbf{Hourly Cost (Spot)} & \textbf{For 2hr Test} \\
\midrule
T4 GPU & \textasciitilde\$0.35 & \$0.70 \\
n1-standard-4 & \textasciitilde\$0.05 & \$0.10 \\
GKE Management & \textasciitilde\$0.10 & \$0.20 \\
\midrule
\textbf{Total} & & \textbf{\textasciitilde\$1-2} \\
\bottomrule
\end{tabular}
\caption{Estimated testing costs using Spot instances}
\end{table}

\begin{tcolorbox}[colback=yellow!10,colframe=orange!50!black,title=\textbf{Important: Cleanup}]
After testing is complete, I will delete the GKE cluster to stop all charges. 

If you create the cluster yourself, please remember to delete it after we finish testing:
\begin{lstlisting}[language=bash]
gcloud container clusters delete piqc-test \
  --zone=us-central1-a --quiet
\end{lstlisting}
\end{tcolorbox}

%==============================================================================
\section{Summary of What I Need}
%==============================================================================

Please choose one of these options:

\begin{enumerate}[label=\textbf{\arabic*.}]
    \item \textbf{Grant me Editor access} to the project (recommended, quickest)
    \item \textbf{Grant minimal roles} (Kubernetes Engine Admin, Compute Admin, Service Usage Admin)
    \item \textbf{Set up the cluster yourself} and share the kubeconfig file with me
\end{enumerate}

\vspace{1em}

\begin{tcolorbox}[colback=blue!5,colframe=blue!50!black,title=\textbf{Information I Need From You}]
\begin{itemize}
    \item \textbf{GCP Project ID:} \_\_\_\_\_\_\_\_\_\_\_\_\_\_\_\_\_\_\_\_\_\_
    \item \textbf{Preferred Option:} 1 / 2 / 3 (circle one)
    \item \textbf{Any specific zone preference?} (default: us-central1-a)
\end{itemize}
\end{tcolorbox}

\vspace{2em}

Thank you for your assistance. Please let me know if you have any questions.

\vspace{2em}
\noindent\textbf{Best regards,}\\
\textit{[Your Name]}

\end{document}
